\section{Глава~\ref{sec:sleep}. Сон}

Режимы сна, повторы и уменьшение синапсов измерены записями нейронов, микродиализом и электронной микроскопией; расписание сна и медленные качели описаны моделями книги, а назначение сна --- в основном гипотезы.

{\footnotesize
\setlength{\tabcolsep}{3pt}\renewcommand{\arraystretch}{1.15}
\begin{longtable}{>{\raggedright}p{0.5cm}>{\raggedright}p{4.0cm}>{\raggedright}p{5.6cm}>{\raggedright}p{2.3cm}>{\raggedright\arraybackslash}p{1.7cm}}
\toprule
№ & Утверждение & Опыт и что измерено & Источник & Статус \\
\midrule\endfirsthead
\toprule
№ & Утверждение & Опыт и что измерено & Источник & Статус \\
\midrule\endhead
1 & Во сне нейроны коры не молчат; меняется режим работы & Долгие записи нейронов коры у крыс: в медленном сне частота остаётся ненулевой, периоды работы чередуются с периодами тишины; после долгого бодрствования частота выше во всех состояниях, после сна --- ниже & \cite{Vyazovskiy2009} & измерено \\
2 & \nt{ACET} высок днём и в быстром сне, низок в медленном (таблица~\ref{tab:sleepmodes}) & Микродиализ в коре у свободно двигающихся кошек: в бодрствовании выброс ацетилхолина на $100$--$175\%$ выше, чем в медленном сне, а в быстром сне --- примерно на уровне спокойного бодрствования & \cite{Marrosu1995} & измерено \\
3 & \nt{NORA} и \nt{SERO} затихают в медленном сне и почти молчат в быстром & Запись отдельных \nt{NORA}-нейронов у крыс и \nt{SERO}-нейронов у кошек по стадиям сна: частота максимальна в бодрствовании, ниже в медленном сне и почти нулевая в быстром & \cite{AstonJonesBloom1981,McGintyHarper1976} & измерено \\
4 & В быстром сне мышцы выключены торможением через \nt{GABA} и \nt{GLYC} & У крыс мерили активность \nt{ACET}-нейронов жевательной мышцы и подавали блокаторы прямо к ним: паралич снимается, только если заблокированы и \nt{GABA}-рецепторы обоих типов, и \nt{GLYC}-рецепторы & \cite{BrooksPeever2012} & измерено \\
5 & Давление сна $S$ --- конденсатор с двумя порогами~\eqref{eq:sleeppressure}, $\tau_r=18.2$~ч, $\tau_d=4.2$~ч & Модель двух процессов; постоянные времени получены из того, как мощность медленных волн ЭЭГ растёт с бодрствованием и падает во сне, а форма порогов --- из времени спонтанного пробуждения & \cite{Borbely1982,Daan1984} & теория \\
6 & Машина сама выходит на $16.1$~ч бодрствования и $8.0$~ч сна (рис.~\ref{fig:twoprocess}) & Расчёт по~\eqref{eq:sleeppressure} с качающимися порогами, скрипт \texttt{models/sleep\_models.py}: $16.1$~ч бодрствования и $7.9$--$8.0$~ч сна & --- & модель \\
7 & После $40$~ч без сна $S=0.90$, но сон длится лишь $8.8$~ч: долг возвращается глубиной, а не длительностью & Модель: \texttt{models/sleep\_models.py} даёт $S=0.90$ и $8.8$~ч. Опыт: после $40.5$~ч без сна у восьми человек в первую восстановительную ночь мощность медленных волн ЭЭГ значимо выше обычной & \cite{Borbely1981} & модель \\
8 & Физически $S$ --- аденозин & Микродиализ у кошек: внеклеточный аденозин в одной из областей, управляющих бодрствованием, растёт за время бодрствования, растёт дальше при лишении сна и падает во время сна. Что $S$ --- только аденозин, не показано & \cite{PorkkaHeiskanen1997,Dunwiddie2001} & гипотеза \\
9 & В медленном сне кора качается: работа несколько сотен мс --- тишина, реже раза в секунду ($<1$~Гц, под наркозом чаще $0.3$--$0.4$~Гц) & Внутриклеточная запись у кошек под наркозом: деполяризация $0.8$--$1.5$~с и долгая гиперполяризация, ритм $<1$~Гц, чаще всего $0.3$--$0.4$~Гц; это же чередование работы и тишины видно в естественном сне (строка~1) & \cite{Steriade1993,Vyazovskiy2009} & измерено \\
10 & Качели --- группа с обратной связью и усталостью~\eqref{eq:updown}; при $b=5$ период $1.1$~с (значение модели) и работа около $35\%$ времени, при $b=1$ работа ровная (рис.~\ref{fig:updown}) & Расчёт, скрипт \texttt{models/sleep\_models.py}: период $1.11$~с, доля времени работы $0.35$ (среднее по целому числу периодов); при $b=1$ доля работы $1.0$. Период модели короче измеренного (строка~9) & --- & модель \\
11 & \nt{ACET} выключает качели и переводит кору в ровную работу & Стимуляция ядра \nt{ACET}-нейронов у кошки блокировала медленный ритм в $79\%$ нейронов коры и переводила их в ровную работу, в основном за счёт исчезновения долгих гиперполяризаций. Ацетилхолин подавляет медленные калиевые токи, создающие усталость & \cite{SteriadeAmzica1993,McCormick1992} & измерено \\
12 & Вспышки ритма $7$--$15$~Гц порождает петля \nt{GLUT}--\nt{GABA} между корой и релейным узлом & В срезах релейного узла зрения у хорька вспышки создаются ответными пачками релейных клеток на торможение от соседнего \nt{GABA}-ядра, которое они сами возбуждают & \cite{vonKrosigk1993,SteriadeMcCormickSejnowski1993} & измерено \\
13 & Повторы дня идут с перемоткой: в гиппокампе примерно в $20$ раз быстрее, в коре --- в $6$--$7$ раз & В медленном сне повторяющиеся последовательности нейронов гиппокампа идут сжатыми во времени. После бега по дорожке последовательности нейронов места повторялись в том же порядке короткими вспышками по $\sim100$~мс, сжатыми примерно в $20$ раз; в коре сжатие в $6$--$7$ раз & \cite{Nadasdy1999,LeeWilson2002,Euston2007} & измерено \\
14 & Усиление согласованности повторов с корой улучшает память (причинная связь) & У крыс после обучения стимуляцией, привязанной к повторам, усиливали их связь с медленными волнами и вспышками ритма в коре: на следующий день память была высокой, у контрольных --- на уровне случая & \cite{Maingret2016} & измерено \\
15 & Назначение повторов --- перемешанное обучение медленной памяти & Теория двух систем обучения и расчёты на искусственных сетях: последовательное обучение портит старое, перемешанное --- нет. Прямого опыта на мозге, что повторы нужны именно для этого, нет & \cite{McClelland1995} & гипотеза \\
16 & После сна синапсы уменьшаются пропорционально размеру, крупные почти не меняются & Трёхмерная электронная микроскопия $6920$ синапсов коры мыши: площадь контакта после сна на $\sim18\%$ меньше, чем после бодрствования, уменьшение пропорционально размеру, крупные устойчивые синапсы не меняются & \cite{deVivo2017} & измерено \\
17 & Главная задача сна --- перенормировка весов & Гипотеза синаптического гомеостаза; строки 1 и 16 с ней согласуются, но не доказывают, что это главная функция & \cite{TononiCirelli2014} & гипотеза \\
18 & Быстрый сон --- испытание модели мира на стенде & Режим (таблица~\ref{tab:sleepmodes}) измерен; что сеть прогоняет модель мира, --- теоретическое предположение, опыта нет & \cite{Hobson2009} & гипотеза \\
19 & Промывка мозга во сне: вопрос открыт & Один опыт: во сне межклеточное пространство на $60\%$ шире и уборка быстрее. Другой, с иным способом измерения: во сне и под наркозом уборка заметно медленнее. Результаты противоречат друг другу & \cite{Xie2013,Miao2024} & гипотеза \\
20 & Местный сон: участки коры бодрствующего животного ненадолго проваливаются в тишину & У крыс после долгого бодрствования нейроны отдельного участка коры коротко замолкают, как в медленном сне, с медленной волной в местной ЭЭГ; в такие моменты крыса чаще ошибается в задаче & \cite{Vyazovskiy2011} & измерено \\
\bottomrule
\end{longtable}}
